% =============================================================
% Abstract
% =============================================================

\begin{abstract}
The rapid growth of large language model (LLM) inference workloads
imposes serving-cost and energy pressures that single-axis
throughput-oriented schedulers cannot address. We study
\emph{cloud-edge collaborative LLM inference scheduling} under five
AIGC-specific physical constraints: model-weight residency
(7--70~GB), KV-cache memory growth, continuous batching, the
memory-bandwidth floor of decode, and the two-phase
prefill/decode request lifecycle. We design an AIGC-aware
proximal policy optimization (PPO) scheduler that exposes these
physics to the policy through state features and reward shaping,
and we compare it against eight strong baselines spanning
heuristic (HEFT, GA, PSO), load-balancing (LeastLoaded,
ShortestQueue), and learned (A3C-R2N2, GNN) approaches.
Across 30 configurations sweeping topology
(edge $\in \{3,5,7\}$), load ($\lambda \in \{0.5,1,2,4,8\}$
req/s), and workload composition (uniform versus large-heavy
model mix) with $N{=}30$ runs each, our scheduler achieves
strict Pareto dominance in the canonical operating regime
($\lambda{=}2$ req/s, uniform mix, edge${=}5$): $+28\%$ \slo{}
attainment and $-7\%$ energy per token relative to the strongest
baseline, with significance at $p<0.05$ on \slo{} against all four
strong comparators and on energy efficiency against three of four
(PSO $p=0.11$, others $p<0.05$). A 12-component ablation,
however, reveals that the joint gain is a holistic system
contribution: only continuous-batching simulation and adequate
pretraining are individually significant; every AIGC reward and
state component is statistically indistinguishable from the full
model when ablated alone. We release the simulator and the 30+
benchmark manifests to support future AIGC scheduling research.
\end{abstract}
